\bookchapter{פרק 12: בידור אודיו-ויזואלי וקריאה}

\chapterintro{אפליקציות המוזיקה כיום דומות למדי בפונקציונליות, אבל לכל אחת יש דגשים משלה – לבחור נכון חוסך הרבה טרחה:}

\booksection{NetEase Cloud Music}

:היתרון שלו הוא קהילת הפלייליסטים ואזור התגובות. אם אתה רוצה לשמוע “שירים שמתאימים ליום גשום” או “מוזיקה אינסטרומנטלית טהורה להירגעות אחרי לימודים”, פשוט תחפש ויש פלייליסטים שמישהו כבר ארגן. אזור התגובות לפעמים יותר מעניין מהשירים. החיסרון הוא שחלק מזכויות השירים אינן מלאות, ויש שירים שהופכים לאפורים תוך כדי האזנה.

\begin{itemize}
  \item \textbf{NetEase Music}:זכויות היוצרים המלאות ביותר, זמרים פופולריים כמו ג'ואו ג'יילון ומיידיי יש בעיקר רק ב-QQ Music. אפשרויות איכות השמע גם עשירות מאוד, חברים יכולים לפתוח איכות שמע ללא אובדן. אם חשוב לך "למצוא את כל השירים שאתה רוצה לשמוע", QQ Music היא הבחירה הבטוחה ביותר.
  \item \textbf{QQ Music}:הבחירה המועדפת על משתמשי אפל. ממשק נקי, ללא פרסומות, איכות שמע טובה, ותומך ב-Lossless וב-Dolby Atmos. אבל אלגוריתם ההמלצות שלו לא מבין את העדפות המשתמשים הסיניים כמו אפליקציות מקומיות, ותרבות הפלייליסטים לא מפותחת כמו NetEase Cloud. מתאים לאנשים שאוהבים חוויה פשוטה ופחות מתעניינים בפונקציות חברתיות.
  \item \textbf{Spotify}:בעל מספר המשתמשים הגדול בעולם, אלגוריתם ההמלצות נחשב לחזק ביותר. "Discover Weekly" ו-"Daily Mix" בעלי דיוק מדהים, ככל שמשתמשים יותר כך הוא מבין אותך יותר. החיסרון הוא שהגרסה החינמית מלאה בפרסומות, וחלק מזכויות השירים הסיניים אינן מלאות. דורש רשת מיוחדת לשימוש.
  \item \textbf{כמה טיפים להאזנה למוזיקה}:אם אתה משתמש באוזניות בלוטות', לך להגדרות אפליקציית המוזיקה ובדוק אם יש "התאמת אפקטים קוליים" או "אקולייזר", כוונן מעט את הבס והטרבל, השיפור מורגש מאוד. בנוסף, האזנה לסטרימינג צורכת הרבה נפח גלישה, זכור להפעיל "נגן Lossless רק ב-WiFi" בהגדרות, או הורד מראש פלייליסטים אהובים להאזנה לא מקוונת.
\end{itemize}

\textbf{2.פודקאסטים:ה“בן לוויה לאוזניים” בנסיעות}פודקאסטים, אנשים שלא שמעו חושבים שזה לא רלוונטי אליהם, ברגע שמתחילים להאזין אי אפשר לחזור אחורה. בעצם זו “גרסת Netflix של תוכניות רדיו” — יש אנשים שמדברים על טכנולוגיה, אחרים על היסטוריה, אחרים קוראים סיפורים, יש כל סוג שתרצה.

\booksection{אפליקציית XiaoYu}

:אפליקציית הפודקאסטים הטובה ביותר בסין, ממשק נקי, ללא פונקציות מבולגנות. “ההמלצה היומית” ו“בחירת העורך” שלה עוזרים לך לגלות תוכניות חדשות, אזור התגובות גם פעיל מאוד. אני אישית מאזין הכי הרבה ל“ימינה ושמאלה”(קטגוריית תרבות והיסטוריה) ו“בלתי ראוי”(קטגוריית נושאים חברתיים), כל פרק מרגיש כמו להאזין לכמה אנשים חכמים משוחחים.

\begin{itemize}
  \item \textbf{פודקאסטים של אפל}:מובנה ב-iPhone, אין צורך להוריד אפליקציה נוספת. האינטגרציה עם CarPlay היא הטובה ביותר, נוח מאוד בנהיגה.
  \item \textbf{טיפים להאזנה}:רוב אפליקציות הפודקאסטים תומכות בהאזנה במהירות 1.2x או 1.5x, אחרי שמתרגלים חוסכים הרבה זמן. אני בדרך כלל מאזין ב-1.2x לסינית, ו-1.5x לאנגלית. בנוסף חשוב להפעיל "הורדה אוטומטית של תוכניות חדשות" ו"מחיקה אוטומטית אחרי השמעה", כך הטלפון שלך לא יתמלא בקבצי שמע שכבר האזנת להם.
  \item \textbf{3.צפייה בסרטונים:הדרך הנכונה מנגנים לסטרימינג}הרבה אנשים לא יודעים, הפורמטים שניתן לנגן בטלפון הם הרבה יותר מאלה שזמינים באפליקציות הסטרימינג. נגן מקומי טוב יכול לפתור הרבה בעיות.
\end{itemize}

\booksection{נגן VLC}

:קוד פתוח וחינמי, כמעט כל פורמט ניתן לנגן — MKV, AVI, FLAC, ואפילו תמונות ISO. אם הורדת סרטון שעובד במחשב אבל לא בטלפון, סביר להניח שזו בעיית פורמט, התקן VLC וזה נפתר. הוא גם תומך בכתוביות חיצוניות, שים קובץ כתוביות .srt באותה תיקייה עם הסרטון, VLC יטען אותן אוטומטית.

\begin{itemize}
  \item \textbf{nPlayer}:אם אתה רוצה חוויה מעודנת יותר, אפליקציה בתשלום זו שווה קנייה. ממשק יפה, מחוות שליטה נוחות — החלק למעלה/למטה בצד שמאל להתאמת בהירות, צד ימין להתאמת עוצמת קול, החלק שמאלה/ימינה להרצה קדימה/אחורה. תומך ב-DTS וב-Dolby Audio, תחושת הנוכחות חזקה הרבה יותר בצפייה בסרטים גדולים.
  \item \textbf{Bilibili (Bilibili)}:כבר לא רק אתר אנימה, עכשיו משאבי הלמידה שם רבים בטירוף. מבחני כניסה, תכנות, צילום, בישול, יש כל מיני מדריכים. תרבות ה"דנמו" (הערות צפות) היא הנשמה של Bilibili, כשצופים בווידאו עם דנמו, מרגיש כאילו יש קבוצת אנשים שצופה איתך. אבל כשצופים בסרטוני למידה מומלץ לכבות דנמו, אחרת קל לאבד ריכוז.
  \item \textbf{שיקוף לטלוויזיה}:מסך הטלפון, גדול ככל שיהיה, לא משתווה לנוחות הצפייה בטלוויזיה. באותה רשת WiFi, פתח את כפתור ה-"TV" בפינה העליונה של אפליקציית הווידאו, ותוכל להקרין את התמונה לטלוויזיה. רוב הטלוויזיות החכמות והמקרנים תומכים בכך, חוויית הצפייה בסדרות וסרטים הרבה יותר טובה. משתמשי iPhone יכולים גם להשתמש ב-AirPlay לשיקוף ישיר של כל המסך.
  \item \textbf{4.קריאת ספרים אלקטרוניים:אפשר לקרוא ספרים היטב גם בטלפון}יש אומרים שקריאה בטלפון פוגעת בעיניים ומקשה על הריכוז, ויש בזה אמת. אבל היתרונות של קריאה בטלפון ברורים — בכל זמן ובכל מקום, ניצול מקסימלי של זמן מפוצל. אני קראתי לא מעט ספרים בטלפון ברכבת התחתית, בתור ולפני השינה.
\end{itemize}

\booksection{WeChat Reading}

:אפליקציית הקריאה עם החוויה הטובה ביותר בסין כרגע. ספרייה מלאה מאוד, יש גם ספרים מודפסים וגם ספרות אינטרנט. היא משתמשת במודל “זמן קריאה תמורת מטבעות ספרים”, ככל שקוראים יותר מקבלים יותר ספרים בחינם, ידידותי מאוד לסטודנטים. פונקציית ה“הדגשה” מאפשרת לראות אילו קטעים אחרים סימנו, לפעמים לראות משפטים טובים שאחרים הדגישו באמת יוצר הזדהות. אבל “טבלת דירוג הקריאה” היא קצת תחרותית, קל להפוך ל“לקרוא כדי להשיג זמן”, מומלץ לכבות.

\begin{itemize}
  \item \textbf{אפליקציית Kindle}:אם יש לך מכשיר Kindle, אפליקציית Kindle בטלפון יכולה לסנכרן את התקדמות הקריאה. איכות העימוד של ספרים אלקטרוניים מקוריים ב-Kindle הרבה יותר טובה מ-PDF פיראטי, אבל חנות הספרים של Kindle בסין כבר נסגרה, צריך לעבור לאזור בינלאומי.
  \item \textbf{Apple Books (ספרים)}:מובנה ב-iPhone, מתאים לקריאת EPUB ו-PDF. ממשק מינימליסטי, אנימציית דפדוף נעימה. אם הורדת קבצי ספרים אלקטרוניים מהאינטרנט, פשוט "שתף לספרים" וזה מיובא.
  \item \textbf{הצעות להגדרות קריאה}:לא משנה באיזו אפליקציה, חשוב להגדיר רקע בגוון צהוב חם (מצב הגנת עיניים), להגדיל מעט את הגופן ולהרחיב את מרווח השורות. לקריאה בלילה פתח "מצב כהה" — טקסט לבן על רקע שחור, הרבה יותר נוח לעיניים. בשום אופן אל תקרא זמן רב באור עמום עם טקסט שחור על רקע לבן בגופן קטן, עייפות העיניים היא מצטברת.
  \item \textbf{5.משחקי טלפון:לשחק בכיף, בלי לפגוע בטלפון}הגישה שלי למשחקי טלפון היא:אפשר לשחק, אבל אל תיתן להם לשלוט בך בחזרה. שימוש מושכל בכלי משחק, משחקים יותר בנוחות, ולא מפספסים דברים חשובים.
\end{itemize}

\booksection{מצב משחק/מרחב משחק}

:לרוב הטלפונים המודרניים יש פונקציה זו. אחרי שמפעילים, הטלפון מנקה אוטומטית אפליקציות רקע, חוסם התראות ומונע נגיעות בשוגג — אתה באמצע קרב קבוצתי, ופתאום שיחה קופצת, מצב משחק יכול לחסום זאת. חלק מהטלפונים תומכים גם במצב ביצועים/E-sports, שמכריח את המעבד לעבוד במלוא הכוח, אבל במחיר של התחממות וצריכת סוללה מהירה יותר.

\begin{itemize}
  \item \textbf{איך לכוונן איכות תמונה וקצב פריימים}:אל תעלה הכל למקסימום בלי מחשבה. קודם ודא שקצב הפריימים (חלקות) טוב, ואז שפר את איכות התמונה. 60fps הרבה יותר נוח מ-30fps, אבל 120fps מכביד מאוד על הטלפון, טלפונים בינוניים לא צריכים להתאמץ. אם אתה משחק במשחקים תחרותיים (כמו Honor of Kings, PUBG), עדיפות קצב הפריימים גבוהה בהרבה מאיכות התמונה — שליטה חלקה הרבה יותר חשובה מדשא יפה.
  \item \textbf{מה לעשות כשהטלפון מתחמם במשחק}:זה נורמלי, אבל חום גבוה לאורך זמן פוגע בחיי הסוללה. כמה שיטות קירור:הסר את הכיסוי לפיזור חום; הורד את בהירות המסך לכ-70\% (בהירות היא מקור חום גבוה); אל תשחק תוך כדי טעינה (הצטברות חום, הסוללה מתה מהר). אם הטלפון חם מכדי להחזיק, צא מיד ותן לו לנוח קצת.
  \item \textbf{כלים למניעת התמכרות}:"זמן מסך" ב-iOS ו"רווחה דיגיטלית" באנדרואיד יכולים להגדיר מגבלת זמן יומית לאפליקציות משחק. אני הגדרתי לעצמי שעה מקסימום בימי עבודה, שעתיים בסופ"ש. כשהזמן נגמר אייקון האפליקציה הופך לאפור ואי אפשר להיכנס. המגבלה הזו באמת מועילה, לא סתם הצגה.
  \item \textbf{משחקי ענן}:אם ביצועי הטלפון שלך לא מספיקים אבל אתה רוצה לשחק במשחקים כבדים, אתה יכול לנסות פלטפורמות משחקי ענן (כמו Tencent START, NetEase Cloud Gaming). המשחק רץ על שרת מרוחק, התמונה מועברת כזרם וידאו לטלפון שלך. דרישות רשת גבוהות אבל כמעט ללא דרישות לביצועי הטלפון. החיסרון הוא שהשהייה בלתי נמנעת, מתאים למשחקי אסטרטגיה ו-RPG, לא מתאים למשחקי FPS.
  \item \textbf{בוא נתנסה:בנה את מערכת הבידור שלך בטלפון}הורד את 3 האלבומים או הפלייליסטים האהובים עליך באפליקציית המוזיקה, וודא שניתן להאזין לא מקוון.
\end{itemize}

\begin{tcbpractice}{לך ל-XiaoYu או פודקאסטים של אפל, הירשם ל-3 ערוצים מעניינים, הורד פרק אחד להאזנה.}
\begin{enumerate}
  \item התקן נגן VLC, נסה לנגן קובץ וידאו שהעתקת מהמחשב.
  \item חפשו ספר שאתם רוצים לקרוא ב-WeChat Reading, הגדירו רקע בגוון חם ומצב הגנת עיניים, כווננו את הגופן לגודל שמתאים לכם.
  \item אם אתה משחק, פתח מצב משחק וראה אילו הגדרות ניתן לייעל (חסימת הפרעות, עדיפות קצב פריימים וכו').
  \item פרק 13:המלצות לאפליקציות כלים
  \item פרק זה בוחר את אפליקציות הכלים שבאמת יכולות לשפר יעילות, לקדם בריאות ולהקל על נסיעות.
\end{enumerate}

\end{tcbpractice}
